```latex
\subsection{The mixed term as a quarter-turn defect}

Recall that
\[
\rho(\theta,\psi)=\rho^+(\theta,\psi)+\rho^-(\theta,\psi)
=
\frac{e^{u(\theta,\psi)}}{I^+}
+
\frac{e^{-u(\theta,\psi)}}{I^-},
\]
and define
\[
\Xi(\theta,\psi)
:=
\frac{1}{2}D(\theta,\psi)^2
+
\frac{1}{2}D\!\left(\theta,\psi+\frac{\pi}{2}\right)^2
-
C(\theta,\psi)^2.
\]
Let
\[
I:=\int_{S^2} F_A(\theta,\psi)\,F_A(\theta,\psi+\pi)\,d(\mu^++\mu^-).
\]
All functions are understood as \(2\pi\)-periodic in the \(\psi\)-variable. 
All pointwise identities below are used only where the displayed quantities are defined; this is sufficient for the integral identity, since any exceptional set has surface measure zero.

\medskip
\noindent\textbf{(a) Quarter-turn invariance of \(\rho\).}
By the \(C_4\)-symmetry of the full cube,
\[
u\!\left(\theta,\psi+\frac{\pi}{2}\right)=u(\theta,\psi).
\]
Since \(I^\pm\) are constants, it follows that
\[
\rho\!\left(\theta,\psi+\frac{\pi}{2}\right)
=
\frac{e^{u(\theta,\psi+\frac{\pi}{2})}}{I^+}
+
\frac{e^{-u(\theta,\psi+\frac{\pi}{2})}}{I^-}
=
\frac{e^{u(\theta,\psi)}}{I^+}
+
\frac{e^{-u(\theta,\psi)}}{I^-}
=
\rho(\theta,\psi).
\]

\medskip
\noindent\textbf{(b) Quarter-turn invariance of \(C\).}
Because \(F_A=\partial_\psi G_A\) and, by the same definition, \(F_B=\partial_\psi G_B\), we have
\[
u(\theta,\psi)=2\pi\bigl(G_A(\theta,\psi)+G_B(\theta,\psi)\bigr),
\]
hence
\[
\partial_\psi u
=
2\pi\bigl(\partial_\psi G_A+\partial_\psi G_B\bigr)
=
2\pi(F_A+F_B)
=
2\pi C.
\]
Differentiating the identity
\[
u\!\left(\theta,\psi+\frac{\pi}{2}\right)=u(\theta,\psi)
\]
with respect to \(\psi\) yields
\[
\partial_\psi u\!\left(\theta,\psi+\frac{\pi}{2}\right)
=
\partial_\psi u(\theta,\psi),
\]
and therefore
\[
C\!\left(\theta,\psi+\frac{\pi}{2}\right)
=
\frac{1}{2\pi}\partial_\psi u\!\left(\theta,\psi+\frac{\pi}{2}\right)
=
\frac{1}{2\pi}\partial_\psi u(\theta,\psi)
=
C(\theta,\psi).
\]

\medskip
\noindent\textbf{(c) Folding the \(\psi\)-integral.}
Since \(d(\mu^++\mu^-)=\rho(\theta,\psi)\sin\theta\,d\psi\,d\theta\) in spherical coordinates and
\[
F_B(\theta,\psi)=F_A(\theta,\psi+\pi),
\]
we can rewrite \(I\) as
\[
I
=
\int_0^\pi\int_0^{2\pi}
F_A(\theta,\psi)\,F_B(\theta,\psi)\,\rho(\theta,\psi)\sin\theta\,d\psi\,d\theta.
\]
Using
\[
F_A=\frac{C+D}{2},
\qquad
F_B=\frac{C-D}{2},
\]
we obtain
\[
F_A F_B=\frac{1}{4}(C^2-D^2),
\]
and therefore
\[
I
=
\frac{1}{4}\int_0^\pi
\left(
\int_0^{2\pi}
\bigl(C(\theta,\psi)^2-D(\theta,\psi)^2\bigr)\rho(\theta,\psi)\,d\psi
\right)
\sin\theta\,d\theta.
\]

Now fix \(\theta\in(0,\pi)\), and for the next few lines suppress the \(\theta\)-dependence from the notation. 
From
\[
D(\psi)=F_A(\psi)-F_A(\psi+\pi)
\]
and the \(2\pi\)-periodicity of \(F_A\), we get
\[
D(\psi+\pi)
=
F_A(\psi+\pi)-F_A(\psi+2\pi)
=
F_A(\psi+\pi)-F_A(\psi)
=
-D(\psi).
\]
Hence also
\[
D\!\left(\psi+\frac{3\pi}{2}\right)
=
D\!\left(\psi+\frac{\pi}{2}+\pi\right)
=
-D\!\left(\psi+\frac{\pi}{2}\right).
\]

We now split the interval \([0,2\pi)\) into the four quarter-turn sectors. 
Using parts \textbf{(a)} and \textbf{(b)},
\[
\int_0^{2\pi} C(\psi)^2\rho(\psi)\,d\psi
=
\sum_{k=0}^{3}
\int_0^{\pi/2}
C\!\left(\psi+\frac{k\pi}{2}\right)^2
\rho\!\left(\psi+\frac{k\pi}{2}\right)\,d\psi
=
4\int_0^{\pi/2} C(\psi)^2\rho(\psi)\,d\psi.
\]
Similarly, using part \textbf{(a)},
\[
\int_0^{2\pi} D(\psi)^2\rho(\psi)\,d\psi
=
\sum_{k=0}^{3}
\int_0^{\pi/2}
D\!\left(\psi+\frac{k\pi}{2}\right)^2
\rho\!\left(\psi+\frac{k\pi}{2}\right)\,d\psi
\]
\[
=
\int_0^{\pi/2}
\left[
D(\psi)^2
+
D\!\left(\psi+\frac{\pi}{2}\right)^2
+
D(\psi+\pi)^2
+
D\!\left(\psi+\frac{3\pi}{2}\right)^2
\right]
\rho(\psi)\,d\psi
\]
\[
=
2\int_0^{\pi/2}
\left[
D(\psi)^2
+
D\!\left(\psi+\frac{\pi}{2}\right)^2
\right]
\rho(\psi)\,d\psi.
\]

Substituting these two folded formulas back into the expression for \(I\), and then restoring the \(\theta\)-dependence, we find
\[
I
=
\int_0^\pi\int_0^{\pi/2}
\left[
C(\theta,\psi)^2
-\frac{1}{2}D(\theta,\psi)^2
-\frac{1}{2}D\!\left(\theta,\psi+\frac{\pi}{2}\right)^2
\right]
\rho(\theta,\psi)\sin\theta\,d\psi\,d\theta.
\]
By the definition of \(\Xi\), this is exactly
\[
I
=
-\int_0^\pi\int_0^{\pi/2}
\Xi(\theta,\psi)\,\rho(\theta,\psi)\sin\theta\,d\psi\,d\theta.
\]

This proves \textbf{(a)}, \textbf{(b)}, and \textbf{(c)}.
```